\chapter{Dysentery}

\section*{Definition and Overview}
Dysentery is a serious form of infectious diarrhoea in which the bowel becomes inflamed and the stools contain blood and mucus, accompanied by cramping abdominal pain and fever. It is caused by bacteria or parasites that invade the lining of the large intestine. The two main types are bacillary dysentery, caused by Shigella species, and amoebic dysentery, caused by the parasite Entamoeba histolytica. Dysentery is a significant cause of illness in developing countries and among travellers, and requires prompt medical treatment because of the risk of dehydration and complications.

\section*{Epidemiology}
Dysentery is most common in regions with poor sanitation and limited access to clean water, where it causes a large burden of illness and death, particularly in young children. Bacillary dysentery occurs worldwide and can cause outbreaks in institutions, child care settings, and among travellers. Amoebic dysentery is more common in tropical and subtropical regions. In Australia, dysentery is seen mainly in returned travellers, in outbreaks, and in people exposed to poor sanitation.

\section*{Aetiology and Pathophysiology}
Dysentery results from infection of the large bowel by invasive organisms.
\begin{itemize}
    \item \textbf{Bacillary dysentery:} Shigella species, and sometimes Salmonella, Campylobacter, and Escherichia coli strains, invade the lining of the colon and cause inflammation and ulceration
    \item \textbf{Amoebic dysentery:} Entamoeba histolytica, a parasite acquired by ingesting contaminated food or water, invades the bowel wall
    \item \textbf{Transmission:} person-to-person spread through poor hand hygiene, and ingestion of contaminated food or water
    \item \textbf{Mechanism:} invasion of the intestinal lining causes bleeding, mucus production, and inflammation, producing bloody diarrhoea
    \item \textbf{Complications of invasion:} in amoebiasis, the parasite can spread through the bloodstream to the liver, causing abscesses
\end{itemize}

\section*{Clinical Features}
Symptoms typically begin one to three days after exposure and can be severe.
\begin{itemize}
    \item Frequent, loose stools containing blood and mucus
    \item Abdominal cramping and pain, with tenderness over the bowel
    \item Fever, malaise, and chills
    \item Urgency to pass stools and a sensation of incomplete emptying
    \item Nausea and sometimes vomiting
    \item Dehydration, with dry mouth, reduced urine, and weakness in severe cases
    \item Amoebic dysentery may develop more gradually, with milder symptoms and less fever
\end{itemize}

\section*{Differential Diagnosis}
Bloody diarrhoea has a range of causes that must be distinguished.
\begin{itemize}
    \item Inflammatory bowel disease, including ulcerative colitis and Crohn's disease
    \item Infectious colitis from other organisms, including Campylobacter and E. coli
    \item Ischaemic colitis in older people
    \item Haemorrhoids and anal fissures, which cause bleeding but not the systemic illness
    \item Diverticulitis
    \item Bowel cancer in older patients
\end{itemize}

\section*{When to Seek Medical Care}
Anyone with bloody diarrhoea, particularly with fever, dehydration, or recent travel, should seek medical care promptly. Young children, older people, pregnant women, and people with weakened immune systems are at higher risk and should be reviewed early.

\textbf{Call 000 (or local emergency services) for:}
\begin{itemize}
    \item Signs of severe dehydration: extreme thirst, dry mouth, no urine for many hours, confusion, or collapse
    \item High fever with severe abdominal pain
    \item Bloody diarrhoea with severe weakness or faintness
    \item Inability to keep down any fluids
    \item Sudden, severe abdominal pain suggesting a serious complication
\end{itemize}

\section*{Diagnosis and Investigations}
Laboratory testing is essential to identify the organism and guide treatment.
\begin{itemize}
    \item \textbf{Stool testing:} microscopy and culture of stool samples identify Shigella, amoeba, and other organisms
    \item \textbf{Blood tests:} full blood count and inflammatory markers; blood cultures in severe illness
    \item \textbf{Sigmoidoscopy:} visualising the inflamed bowel may be needed in atypical cases
    \item \textbf{Imaging:} ultrasound or CT to detect amoebic liver abscess or bowel complications
    \item \textbf{Hydration assessment:} clinical assessment of the degree of dehydration
\end{itemize}

\section*{Management}
Treatment depends on the organism and the severity of illness.
\begin{itemize}
    \item \textbf{Rehydration:} oral rehydration solution for mild to moderate dehydration; intravenous fluids for severe cases
    \item \textbf{Antibiotics:} for bacillary dysentery, to shorten the illness and prevent spread; for amoebic dysentery, antiparasitic medicines
    \item \textbf{Antibiotic stewardship:} in mild shigellosis, some clinicians avoid antibiotics to reduce resistance, reserving them for severe cases
    \item \textbf{Symptom relief:} medicines for fever, though anti-diarrhoeal drugs are generally avoided because they may worsen invasive infection
    \item \textbf{Rest and nutrition:} small, bland meals as tolerated once vomiting settles
    \item \textbf{Public health measures:} notification of cases and advice to prevent spread to others
\end{itemize}

\section*{Complications}
\begin{itemize}
    \item Severe dehydration and electrolyte disturbance, especially in children
    \item Toxic megacolon, a dangerous dilatation of the colon
    \item Perforation of the bowel, causing peritonitis
    \item Septicaemia, particularly in people with weakened immunity
    \item Amoebic liver abscess, which requires specific treatment
    \item Reactive arthritis and other post-infectious complications
\end{itemize}

\section*{Prognosis and Outcomes}
With prompt diagnosis and treatment, most people with dysentery recover fully within one to two weeks. The illness is more dangerous in young children, older people, and those with malnutrition or weakened immunity, but appropriate rehydration and antimicrobial treatment greatly reduce the risk of serious outcomes. Amoebic liver abscess responds well to treatment when detected.

\section*{Prevention and Self-Care}
\begin{itemize}
    \item Wash hands thoroughly with soap and water after the toilet and before handling food
    \item Practise safe food and water hygiene when travelling, including drinking bottled water and eating freshly cooked food
    \item Avoid swimming in water that may be contaminated
    \item Stay home from work, school, or child care while symptoms are present
    \item Drink rehydration fluids early in the illness
    \item See a doctor promptly for bloody diarrhoea rather than self-treating
\end{itemize}

\section*{Sources and Further Reading}
The following reputable sources provide further information for healthcare students and patients seeking reliable, current advice about dysentery.
\begin{itemize}
    \item Healthdirect Australia. \textit{Dysentery and gastrointestinal infections.} https://www.healthdirect.gov.au/
    \item NHS. \textit{Diarrhoea and vomiting: causes and treatment.} https://www.nhs.uk/conditions/diarrhoea-and-vomiting/
    \item Centers for Disease Control and Prevention. \textit{Shigellosis and amoebiasis information.} https://www.cdc.gov/
    \item World Health Organization. \textit{Diarrhoeal disease and dysentery.} https://www.who.int/
    \item Mayo Clinic. \textit{Intestinal infection and diarrhoea.} https://www.mayoclinic.org/diseases-conditions/diarrhea/
    \item MSD Manual. \textit{Dysentery: professional overview.} https://www.msdmanuals.com/
\end{itemize}
